\section{Parameter collisions and sharp totals}\label{sec:bounds}
The sharp bound concerns the whole periodic set, not a maximum
single period. All rows that coexist must be counted.

\begin{proof}[Proof of Theorem~\ref{thm:bound}]
First take $a\notin\{1,-1\}$. Beyond fixed points and two-cycles,
the only possible row is $001$, which requires $a^2=-1$.
The parameters $K_a$ and $-3K_a$ are distinct, because equality
would give $(a+1)^2=0$. If $a^2=-1$, the value $C=K_a$
contributes exactly two fixed points and one six-cycle, totaling
eight. For other such determinants the total is at most two.

For $a=1$, all possibly nonempty rows give
\begin{equation}\label{eq:a-plus}
 (C,\text{point contribution})=(1,2),\ (-3,2),\ (0,4),\ (-1,6).
\end{equation}
These constants are distinct in characteristic different from two,
except for the equality $-3=0$ in characteristic three.
That collision combines a two-cycle and a four-cycle, totaling
six, and does not meet the six-point three-cycle locus. Thus
the maximum for $a=1$ is six in every allowed characteristic.

For $a=-1$, one has $K_a=0$. At $C=0$, two fixed points
and one two-cycle give four points. The word $01$ gives four
points at $C=-1$. Outside characteristic three there are no
further rows for this determinant. In characteristic three,
$00011$ contributes ten additional points at the same $C=-1$,
while $0111$ gives eight points at $C=1$. The constants
$0,-1,1$ are distinct. The total is therefore fourteen
exactly on \eqref{eq:equality}.

These cases exhaust $k^*$. Outside characteristics two and three,
the eight-point case exists exactly when $-1$ has a square root
in $k$. Otherwise $a=1$, $C=-1$ gives six points.
Every asserted maximum is attained by taking $P=t$ over the
stated field. Finally, every row with period greater than two
requires $a=1$, $a=-1$ or $a^2=-1$, all implying $a^4=1$.
\end{proof}

The maximal characteristic-three locus can also be displayed
without an implicit field extension. Fix any nonconstant $P$
and write $+$ and $-$ in front of $P$ literally as its two
field signs. For $a=-1$, $C=-1$ and $\ch k=3$, the coordinate
words in Table~\ref{tab:sharp} are periodic under the recurrence.
Consecutive entries, cyclically paired, are the ordinary points.
The same $P$ is used in all three rows.

\begin{table}[htbp]
\centering
\begin{tabular}{@{}cl@{}}
\toprule
Least period & Cyclic coordinate word\\
\midrule
4 & $(P+1,\ P-1,\ -P+1,\ -P-1)$\\
5 & $(P+1,\ P,\ P,\ P-1,\ -P)$\\
5 & $(-P+1,\ -P,\ -P,\ -P-1,\ P)$\\
\bottomrule
\end{tabular}
\caption{All fourteen points on the equality locus: a word
$(y_0,\ldots,y_{\ell-1})$ represents the $\ell$ ordered pairs
$(y_i,y_{i+1})$, with indices modulo $\ell$. Repeated individual
coordinates do not merge ordered pairs. The words are obtained
from \eqref{eq:word-reconstruction} using $01$ and $00011$;
Proposition~\ref{prop:exact} proves their mutual disjointness.}
\label{tab:sharp}
\end{table}

\begin{corollary}[Ordinary return counts]\label{cor:zeta}
Let $m_j$ be the number of cycles of length $j$ obtained by adding
the applicable rows in Table~\ref{tab:atlas}. Then
\begin{equation}\label{eq:zeta}
 \#\Fix(H_{a,c}^n)(k(t))=\sum_{j\mid n}j m_j,\qquad
 \zeta_{H,k(t)}(z)=\prod_j(1-z^j)^{-m_j}.
\end{equation}
The product is finite; when there are no periodic points it equals one.
\end{corollary}
\begin{proof}
A cycle of least period $j$ contributes all its $j$ points at
time $n$ exactly when $j\mid n$. This proves the first identity.
For the ordinary Artin--Mazur definition, as formal power series
at zero,
\[
 \log\zeta_{H,k(t)}(z)
 =\sum_{n\ge1}\frac{z^n}{n}\sum_{j\mid n}j m_j
 =-\sum_jm_j\log(1-z^j).
\]
Exponentiation gives the second identity.
\end{proof}

For example, the equality locus has
$\zeta_{H,k(t)}(z)=(1-z^4)^{-1}(1-z^5)^{-2}$.
This records the same finite rational periodic set and is not a
second classification problem or a Frobenius zeta of a variety.
